\documentclass[11pt]{article}
\usepackage[margin=1in]{geometry}
\usepackage{hyperref}
\hypersetup{colorlinks=true, linkcolor=blue, urlcolor=blue, citecolor=blue}
\newcommand{\refl}[2]{\href{#2}{[#1]}}
\usepackage{parskip}
\usepackage{titlesec}
\usepackage{authblk}
\titlespacing*{\section}{0pt}{14pt}{5pt}
\titlespacing*{\subsection}{0pt}{11pt}{4pt}

\newcommand{\uAlonso}{https://www.hup.harvard.edu/books/9780674537019}
\newcommand{\uMokh}{https://doi.org/10.1068/a36218}
\newcommand{\uRamani}{https://www.nber.org/papers/w28876}
\newcommand{\uGupta}{https://doi.org/10.1016/j.jfineco.2022.06.003}
\newcommand{\uMond}{https://www.nber.org/papers/w30041}
\newcommand{\uPnas}{https://doi.org/10.1073/pnas.2408930121}
\newcommand{\uHoward}{https://doi.org/10.1016/j.jfineco.2023.103705}
\newcommand{\uLeslie}{https://doi.org/10.1177/0160017610365538}
\newcommand{\uZillow}{https://www.zillow.com/research/data/}
\newcommand{\uAcs}{https://www.census.gov/data/developers/data-sets/acs-5year.html}
\newcommand{\uDingel}{https://doi.org/10.1016/j.jpubeco.2020.104235}
\newcommand{\uGaz}{https://www.census.gov/geographies/reference-files/time-series/geo/gazetteer-files.html}
\newcommand{\uLodes}{https://lehd.ces.census.gov/data/}
\newcommand{\uCall}{https://www.nber.org/papers/w32117}
\newcommand{\uDelm}{https://doi.org/10.1080/00045608.2015.1096188}
\newcommand{\uDelmb}{https://doi.org/10.1111/gean.70006}
\newcommand{\uFons}{https://doi.org/10.1111/jofi.13398}

\title{Project Plan: Remote Work and Urban Housing Prices}
\author[1]{Sai Krishna Varshith Danda}
\author[1]{Gayathri Kota}
\author[1]{Eshitha Chowdary Kaka}
\author[1]{Zain Khan}
\affil[1]{Arizona State University, Tempe, AZ 85281, USA}
\date{September 27, 2026}
\setlength{\parskip}{4pt}

\begin{document}
\maketitle

\section{Background, Motivation, and Research Problem}

The monocentric city model predicts that housing prices fall with distance from
the central business district (CBD), because households farther out pay more to
commute, and that this gradient flattens when commuting becomes cheaper
\refl{1}{\uAlonso}. Remote work acts like a fall in commuting cost: a worker who
goes in two days a week instead of five pays the cost of distance less often.
Researchers raised this possibility decades ago \refl{2}{\uMokh}, but remote work
was too rare to test at the scale of a housing market until 2020.

The pandemic supplied that test. Ramani and Bloom \refl{3}{\uRamani} found
households and businesses moving out of dense city centers toward lower-density
suburbs, a pattern they named the Donut Effect, and Gupta et al.\
\refl{4}{\uGupta} measured the same change through the price gradient, which went
from clearly negative to nearly flat in most U.S.\ metros during 2020. Mondragon
and Wieland \refl{5}{\uMond} addressed cause by using each metro's pre-pandemic
remote-work share as the measure of exposure, so the result could not be driven
by remote workers moving into areas whose prices were already rising. They
estimated that remote work explains more than half of the 18.9 percent rise in
real U.S.\ house prices from 2019 to 2023.

What remains unsettled is persistence. Return-to-office requirements expanded
from 2023, but remote work settled far above its pre-pandemic level rather than
reverting, and Ramani, Alcedo, and Bloom \refl{6}{\uPnas} found the Donut Effect
still present in 2024. Others expect it to fade as builders respond
\refl{7}{\uHoward}. Most published estimates of within-metro gradients stop in
2022 or 2023; those extending further work at the metro level or focus on the
largest, densest cities.

Our research problem follows from that gap. We ask whether the flattening of the
within-metro price gradient persisted through 2025 in high-growth U.S.\ Sun Belt
metropolitan areas, and how neighborhood characteristics relate to the price
paths we observe. Sun Belt metros matter here for a specific reason. Ramani and
Bloom \refl{3}{\uRamani} found sizeable donuts in large cities and essentially
none in small ones, and much of the evidence is measured relative to a single
main center. Many Sun Belt metros grew during the automobile era and have
several employment centers rather than one \refl{8}{\uLeslie}. If distance from
a single downtown was always a weak predictor of prices there, a study that
measures only that distance may understate what remote work did.

\section{Sub-Questions and Hypotheses}

\subsection{Sub-Question 1}

How did the relationship between distance from the central business district and
ZIP-level housing price growth change between 2018 and 2025 in Sun Belt metros?

\textbf{Hypothesis.} We expect the gradient to flatten sharply between 2020 and
2022, consistent with the published record, and then hold rather than reverse
through 2025. We do not expect a return to the 2018--2019 gradient. We also
expect the flattening to be smaller in magnitude than the estimates reported for
large, dense metros, because these metros had a weaker downtown premium to begin
with. A finding that the gradient fully reversed would contradict
\refl{6}{\uPnas} and would point toward return-to-office requirements outweighing
the residual effect of hybrid work.

\subsection{Sub-Question 2}

Do ZIP codes with greater pre-pandemic remote-work potential show different price
growth after 2020 than otherwise comparable ZIP codes?

\textbf{Hypothesis.} We expect a positive relationship: ZIP codes whose residents
were already more able to work from home before 2020 should show stronger price
growth afterward. We expect this relationship to be strongest in 2021--2022 and
to weaken but not disappear through 2025. We measure exposure using data that
ends before 2020, following \refl{5}{\uMond}, so that the measure cannot have
been produced by the price changes we are trying to explain. We expect the
relationship to be weaker in our sample than in metro-level national estimates,
since Sun Belt neighborhoods are less sharply sorted by occupation than those in
the largest coastal metros.

\subsection{Sub-Question 3}

What distinct price trajectories emerge when ZIP codes are grouped by the shape
of their 2018--2025 price paths, and what characteristics do those groups share?

\textbf{Hypothesis.} We expect roughly three to five recognizable groups:
sustained appreciation, sharp appreciation followed by correction after 2022,
delayed appreciation, and relative stagnation. We expect distance from the CBD to
separate these groups less cleanly than it would in a monocentric metro, and we
expect pre-pandemic remote-work potential, housing stock composition, and
proximity to secondary employment centers to do more of the separating. If
distance turns out to be the dominant discriminator after all, that would be
evidence against the polycentric argument in Section 1.

\section{Data}

\textbf{Zillow Home Value Index (ZHVI)} \refl{9}{\uZillow}. Our primary outcome.
Free monthly CSVs at ZIP-code level, covering roughly 26{,}000 ZIP codes. We use
the all-homes, mid-tier, smoothed and seasonally adjusted series, which reports
the typical value for homes in the 35th to 65th percentile range. Because it is
built from Zillow's own valuation models and smoothed, it may understate
short-term volatility; we treat this as a limitation rather than a reason to
prefer a noisier series. Files are published in wide format, one column per
month, and will be reshaped into a ZIP-month panel.

\textbf{Zillow Observed Rent Index (ZORI)} \refl{9}{\uZillow}. A secondary
outcome, and the literature review gave us a specific reason to want it. Mortgage
rates rose sharply in 2022 and affect purchase prices across a whole metro, which
can obscure changes in the gap between central and suburban neighborhoods. Rents
do not depend directly on mortgage financing, so comparing rent trends against
price trends is a useful check on whether a post-2022 pattern reflects interest
rates or location preferences. The constraint is coverage: ZORI reports roughly
8{,}400 ZIP codes against ZHVI's 26{,}000, because Zillow needs sufficient rental
listings to compute it. We will report coverage for our eight metros before
deciding how much weight rent carries.

\textbf{American Community Survey (ACS) 5-year estimates} \refl{10}{\uAcs}.
Retrieved through the Census Data API at ZIP Code Tabulation Area level. We pull
two vintages. The 2015--2019 release is our pre-pandemic baseline, used for
Sub-Question 2 because it predates the price changes we study. A recent release
describes current neighborhood composition. Tables of interest include means of
transportation to work (B08006, which contains a ``worked from home'' category),
occupation by sex (C24010) for an occupation-based measure of remote-work
capacity following \refl{11}{\uDingel}, median household income (B19013), and
units in structure (B25024). ZCTAs approximate postal ZIP codes but do not match
them exactly, so a share of rows will not join; we will quantify and report that
match rate rather than silently dropping it.

\textbf{Census Gazetteer files} \refl{12}{\uGaz}. ZCTA centroid latitude and
longitude, used to compute distance from each ZIP code to its metro's center.
Plain download, no API key.

\textbf{LEHD Origin-Destination Employment Statistics (LODES)}
\refl{13}{\uLodes}. Job counts at census block level, released state by state.
The workplace area characteristics file reports jobs by the block where they are
located, which is what we would use to identify secondary employment centers so
that Sub-Question 3 can test proximity to subcenters rather than distance to a
single downtown. This is the least certain of our sources and may be dropped if
the processing cost proves too high for our timeline.

\textbf{Scope.} Eight metros: Phoenix, Austin, Atlanta, Dallas, Tampa, Charlotte,
Nashville, and Las Vegas. Two pre-pandemic years and roughly five years afterward
give us a baseline and a post-return-to-office window.

\section{Planned Exploratory Data Analysis}

We will begin with data quality rather than findings. The ZIP-to-ZCTA join is the
first thing to characterize: how many ZIP codes match, which ones do not, and
whether non-matches cluster anywhere that would bias the sample. We will report
ZORI coverage per metro in the same pass.

The exploratory work is organized around four things we want to see.

\textbf{Distributions.} Price levels and growth rates across ZIP codes, by metro
and by year, using histograms and box plots. We expect right-skewed price levels
and want to know whether to work in logs.

\textbf{The gradient itself.} Scatter plots of price growth against distance from
the center, one panel per metro, repeated by year as small multiples. This is the
visual form of Sub-Question 1, and we expect the slope to be the thing that
changes across panels.

\textbf{Spatial structure.} ZIP-level choropleth maps of price growth for each
metro, which should make polycentricity visible if it is there. A monocentric
metro should show smooth rings; a polycentric one should not.

\textbf{Relationships.} Pre-pandemic remote-work potential against subsequent price
growth, plotted and binned, plus a correlation matrix across the neighborhood
characteristics we plan to use as controls.

Legibility matters here, since these plots are the main way the analysis gets
communicated: consistent color scales across metros, readable axis labels, and
legends that do not require the reader to guess.

\section{Methodological Approach}

For Sub-Question 1 we will estimate an event-study specification interacting
distance from the center with period indicators, with ZIP and metro-by-month
fixed effects. The metro-by-month terms matter for a reason the literature review
made clear: they compare each ZIP code only against others in the same metro and
month, so metro-wide movements such as the 2022 rate shock do not masquerade as
changes in the gradient. Distance is continuous rather than binary, and Callaway,
Goodman-Bacon, and Sant'Anna \refl{14}{\uCall} show that two-way fixed effects
with continuous treatments admit several interpretations and that comparing
effects across treatment levels requires stronger assumptions than the binary
case. We will follow their guidance on aggregation and report results
accordingly.

Sub-Question 2 extends this with pre-pandemic remote-work potential and
demographic controls. Sub-Question 3 applies clustering to normalized price
trajectories, choosing the number of groups by silhouette width and checking that
the resulting clusters are distinct rather than arbitrary partitions of a
continuum. Trajectory clustering has been used in neighborhood research
\refl{15}{\uDelm}, \refl{16}{\uDelmb} but, as far as we found, not on
post-pandemic price paths.

Two limits shape how we will state our results. The pandemic affected every
neighborhood simultaneously, so no untreated comparison group exists. And
mortgage lock-in has reduced mobility substantially since 2022
\refl{17}{\uFons}, which means a stable gradient could reflect households being
unable to move rather than not wanting to. We will therefore present our findings
as patterns consistent or inconsistent with a remote-work explanation, not as
causal estimates.

\begin{thebibliography}{17}\setlength{\itemsep}{0pt}

\bibitem{alonso} W. Alonso. \textit{Location and Land Use: Toward a General
Theory of Land Rent}. Harvard University Press, 1964.

\bibitem{mokhtarian} P. L. Mokhtarian, G. O. Collantes, and C. Gertz.
Telecommuting, residential location, and commute-distance traveled: Evidence from
State of California employees. \textit{Environment and Planning A},
36(10):1877--1897, 2004. \href{\uMokh}{doi.org/10.1068/a36218}.

\bibitem{ramani} A. Ramani and N. Bloom. The donut effect of COVID-19 on cities.
NBER Working Paper 28876, 2021. \href{\uRamani}{nber.org/papers/w28876}.

\bibitem{gupta} A. Gupta, V. Mittal, J. Peeters, and S. Van Nieuwerburgh.
Flattening the curve: Pandemic-induced revaluation of urban real estate.
\textit{Journal of Financial Economics}, 146(2):594--636, 2022.
\href{\uGupta}{doi.org/10.1016/j.jfineco.2022.06.003}.

\bibitem{mondragon} J. A. Mondragon and J. Wieland. Housing demand and remote
work. NBER Working Paper 30041, 2022, revised May 2025.
\href{\uMond}{nber.org/papers/w30041}.

\bibitem{pnas} A. Ramani, J. Alcedo, and N. Bloom. How working from home reshapes
cities. \textit{Proceedings of the National Academy of Sciences},
121(45):e2408930121, 2024. \href{\uPnas}{doi.org/10.1073/pnas.2408930121}.

\bibitem{howard} G. Howard, J. Liebersohn, and A. Ozimek. The short- and long-run
effects of remote work on U.S. housing markets. \textit{Journal of Financial
Economics}, 150(1):166--184, 2023.
\href{\uHoward}{doi.org/10.1016/j.jfineco.2023.103705}.

\bibitem{leslie} T. F. Leslie. Identification and differentiation of urban
centers in Phoenix through a multi-criteria kernel-density approach.
\textit{International Regional Science Review}, 33(2):205--235, 2010.
\href{\uLeslie}{doi.org/10.1177/0160017610365538}.

\bibitem{zillow} Zillow Group. Housing data (ZHVI and ZORI). Zillow Research,
2026. \href{\uZillow}{zillow.com/research/data}.

\bibitem{acs} U.S. Census Bureau. American Community Survey 5-year data, Census
Data API. \href{\uAcs}{census.gov/data/developers}.

\bibitem{dingel} J. I. Dingel and B. Neiman. How many jobs can be done at home?
\textit{Journal of Public Economics}, 189:104235, 2020.
\href{\uDingel}{doi.org/10.1016/j.jpubeco.2020.104235}.

\bibitem{gazetteer} U.S. Census Bureau. Gazetteer files, ZIP Code Tabulation
Areas. \href{\uGaz}{census.gov/geographies/reference-files}.

\bibitem{lodes} U.S. Census Bureau. LEHD Origin-Destination Employment Statistics
(LODES), version 8. \href{\uLodes}{lehd.ces.census.gov/data}.

\bibitem{callaway} B. Callaway, A. Goodman-Bacon, and P. H. C. Sant'Anna.
Difference-in-differences with a continuous treatment. NBER Working Paper 32117,
2024. \href{\uCall}{nber.org/papers/w32117}.

\bibitem{delmelle} E. C. Delmelle. Mapping the DNA of urban neighborhoods:
Clustering longitudinal sequences of neighborhood socioeconomic change.
\textit{Annals of the American Association of Geographers}, 106(1):36--56, 2016.
\href{\uDelm}{doi.org/10.1080/00045608.2015.1096188}.

\bibitem{delmelle2} E. C. Delmelle, I. Nilsson, and N. Duma. Time series
clustering for exploring neighborhood dynamics. \textit{Geographical Analysis},
57(3):402--420, 2025. \href{\uDelmb}{doi.org/10.1111/gean.70006}.

\bibitem{fonseca} J. Fonseca and L. Liu. Mortgage lock-in, mobility, and labor
reallocation. \textit{Journal of Finance}, 79(6):3729--3772, 2024.
\href{\uFons}{doi.org/10.1111/jofi.13398}.

\end{thebibliography}

\end{document}